\section{Whitney implies Toda}\label{sec:whitney}
In this section, we consider a different presentation of the quantum K ring, named the
{\em quantum~K Whitney presentation}.
This presentation quantizes
 relations $\lambda_y(\cS_i) \cdot \lambda_y(\cS_{i+1}/\cS_i) = \lambda_y(\cS_{i+1})$
 satisfied by the tautological subbundles in $\K_T(\Fl(\bfr,n))$. Informally, the
 Whitney presentation contains more (geometric) information than the Toda presentation,
 as it involves more generators, corresponding to the $\lambda_y$
 classes of the tautological subbundles, and their quotients. In~contrast, the Toda presentation
 only involves the quotient bundles.




The quantization was conjectured in~\cite{gu2023quantum,GMSXZZ:QKW},
generalizing the conjectures from
\cite{Gu:2020zpg} for Grassmannians. These conjectures have been proved in~\cite{gu2022quantum}
for Grassmannians, and in~\cite{GMSXZZ:QKW} for~$\Fl(1,n-1;n)$ case. The general case was recently announced in~\cite{huq2024quantum} using the abelian/non-abelian correspondence.
We note that the results in~\cite{huq2024quantum} are logically independent on those from~\cite{maeno2023presentation}, which were used to obtain the Toda presentation in the previous section.


Our main result of this section is that eliminating the additional variables of the Whitney presentation yields the Toda presentation. As an aside, we note that the proof of Theorem~\ref{thm:main} can be easily modified to show that the quantum K Whitney presentation of $\Fl(\bfr,n)$ follows from that of $\Fl(n)$.
We leave the details of this proof to the reader.

In what follows, $T$ can be a maximal torus in $\GL_n$.
Let
\[
 {X}^{(j)}=\bigl(X^{(j)}_1,\dots,X^{(j)}_{r_j}\bigr) \qquad\text{and}\qquad {Y}^{(j)}=\bigl(Y^{(j)}_1,\dots,
Y^{(j)}_{r_{j+1}-r_j}\bigr)
\]
denote formal variables for $j=1,\dots,k$ and denote by $X^{(k+1)}\coloneqq (T_1,\dots,T_n)$ the equivariant parameters in $\K_T(\pt)$. Let \smash{$e_\ell\bigl({X}^{(j)}\bigr)$} and \smash{$e_\ell\bigl({Y}^{(j)}\bigr)$} be the $\ell$-th elementary symmetric polynomials in ${X}^{(j)}$ and ${Y}^{(j)}$, respectively. Define the ring
 \[
 S=\K_T(\pt)\bigl[e_1\bigl(X^{(j)}\bigr),\dots, e_{r_j}\bigl(X^{(j)}\bigr),e_1\bigl(Y^{(j)}\bigr),\dots,e_{r_{j+1}-r_j}\bigl(Y^{(j)}\bigr)\bigr]_{j=1}^k ,
 \]
 and the ideal $I_Q\subset S[\![Q]\!]=S[\![Q_1,\dots,Q_k]\!]$
generated by the coefficients of $y$ in
\begin{gather}\label{eqn:qkrel}
 \prod_{\ell=1}^{r_j}\bigl(1+y X^{(j)}_\ell\bigr)\prod_{\ell=1}^{r_{j+1}-r_j}\bigl(1+y Y^{(j)}_\ell\bigr)-\prod_{\ell=1}^{r_{j+1}}\bigl(1+y X^{(j+1)}_\ell\bigr)\\
 \quad{}+y^{{r_{j+1}-r_j}}\frac{Q_j}{1-Q_j}\prod_{\ell=1}^{r_{j+1}-r_j}Y^{(j)}_\ell\Biggl(\prod_{\ell=1}^{r_j}\bigl(1+yX^{(j)}_\ell\bigr) -\prod_{\ell=1}^{r_{j-1}}\bigl(1+yX^{(j-1)}_\ell\bigr)\Biggr), \qquad j=1,\dots, k. \nonumber
\end{gather}
It was conjectured in~\cite{gu2023quantum,GMSXZZ:QKW} and proved in
\cite{huq2024quantum} that
there is an isomorphism of $\K_T(\pt)[\![Q]\!]$-algebras
 \begin{equation}\label{E:QKWhitney}
 \Phi\colon\ {S[\![Q]\!]}/I_Q\to \QK_T\br{\Fl(\bfr,n)}
 \end{equation}
 sending
 \[ e_\ell\bigl(X^{(j)}\bigr) \mapsto \wedge^\ell(\cS_{j}) \qquad \textrm{and} \qquad e_\ell\bigl(Y^{(j)}\bigr) \mapsto \wedge^\ell(\cS_{{j+1}}/\cS_{j}) . \]
 We refer to this as the (quantum K) {\em Whitney presentation}.

 \begin{Proposition}\label{prop:whit-toda}
 There is a natural isomorphism
 \[ S[\![Q]\!]/I_Q \simeq R\dbb{Q}/J_Q ,\]
 obtained by eliminating the indeterminates \smash{$X^{(j)}_\ell$}. In particular,
 the Whitney relations from~\eqref{eqn:qkrel} imply the
 Toda relations from~\eqref{eqn:whit-short-gen}.
 \end{Proposition}
 \begin{proof}
 Let
 \begin{align*}
 A_j=\prod_{\ell=1}^{r_{j+1}-r_j}\bigl(1+y Y^{(j)}_\ell\bigr)+B_j,\qquad B_j=y^{{r_{j+1}-r_j}}\frac{Q_j}{1-Q_j}\prod_{\ell=1}^{r_{j+1}-r_j}Y^{(j)}_\ell,
 \end{align*}
 so that~\eqref{eqn:qkrel} becomes
 \begin{align}\label{eqn:recur}
 A_j\prod_{\ell=1}^{r_j}\bigl(1+yX^{(j)}_\ell\bigr) -B_j\prod_{\ell=1}^{r_{j-1}}\bigl(1+yX^{(j-1)}_\ell\bigr) -\prod_{\ell=1}^{r_{j+1}}\bigl(1+y X^{(j+1)}_\ell\bigr).
 \end{align}
 Note that by Lemma~\ref{L:r-det}, relations given by~\eqref{eqn:recur} are equivalent to those given by
 \begin{align*}%\label{eqn:elim}
 \prod_{\ell=1}^{r_{j+1}}\bigl(1+y X^{(j+1)}_\ell\bigr)-
 &\begin{vmatrix}
 A_0 & B_1 & & & \\
 1 & A_1 & B_2 & & \\ & \ddots & \ddots & \ddots & \\ & & 1 & A_{j-1} & B_{j}\\ & & & 1 & A_{j}
 \end{vmatrix}\qquad \text{for $1\leq j\leq k$}.
 \end{align*}
 As a consequence, we can eliminate $e_1\bigl(X^{(j)}\bigr),\dots, e_{r_j}\bigl(X^{(j)}\bigr)$ for $2\leq j\leq k$, and be left with the relation~\eqref{eqn:whit-short-gen}.
 \end{proof}
\begin{Remark}\label{rmk:QKW} We note that our methods from the previous section can be adapted easily to show that $\Phi$ is an isomorphism for all partial flag varieties if and only if it is an isomorphism for $\Fl(n)$. \end{Remark}
